\documentclass[UTF8,fontset=fandol,11pt]{ctexart}
\usepackage[a4paper,margin=25mm]{geometry}
\usepackage{amsmath,amssymb,amsthm,mathtools}
\usepackage{xcolor}
\usepackage{booktabs,array,enumitem,fancyhdr,hyperref}
\hypersetup{colorlinks=true,linkcolor=blue!50!black,urlcolor=blue!50!black}
\setlength{\emergencystretch}{2em}
\setlength{\parskip}{0.3em}
\linespread{1.16}
\allowdisplaybreaks
\newtheorem{theorem}{定理}[section]
\newtheorem{proposition}[theorem]{命题}
\newtheorem{lemma}[theorem]{引理}
\theoremstyle{definition}
\newtheorem{definition}[theorem]{定义}
\theoremstyle{remark}
\newtheorem{remark}[theorem]{说明}
\newcommand{\dd}{\,\mathrm d}
\newcommand{\ii}{\mathrm i}
\newcommand{\Rea}{\operatorname{Re}}
\newcommand{\Ima}{\operatorname{Im}}
\newcommand{\Li}{\operatorname{Li}}
\newcommand{\supp}{\operatorname{supp}}
\pagestyle{fancy}
\fancyhf{}
\fancyhead[L]{\small 学习笔记}
\fancyhead[R]{\small 初稿}
\fancyfoot[C]{\thepage}
\setlength{\headheight}{15pt}
\author{}
\date{2026年10月22日初稿}
\makeatletter
\renewcommand{\maketitle}{%
  \thispagestyle{plain}%
  \begin{center}
    {\LARGE\bfseries \@title\par}
    \vspace{0.6em}
    {\small \@date\par}
  \end{center}
  \vspace{0.6em}}
\makeatother
\newcommand{\draftnotice}{\noindent{\small
本文为供逐篇审阅的学习笔记初稿，不是新增研究成果。
文中区分直接证明与引用结果；正文不依赖本文。
本轮已作公式、条件及排版初步核查，尚待作者逐句审定。}\par\medskip}

\title{积分计数函数的显式公式}
\begin{document}
\maketitle
\draftnotice
\section{积分以后，新的对象是什么}
记 $\Lambda(n)$ 为 von Mangoldt 函数，$\psi(x)=\sum_{n\le x}\Lambda(n)$。
对 $x\ge1$ 定义
\[
 \psi_1(x)=\int_1^x\psi(u)\dd u
 =\sum_{n\le x}\Lambda(n)(x-n).
\]
有限区间只有有限个跳跃，故后一个等式可逐项积分获得。
积分端点处 $\psi$ 取左值、右值或左右平均，不影响积分。
$\psi_1$ 连续，但在素数幂处的导数一般仍有跳跃。
本篇推导其显式公式，并证明零点级数绝对收敛；不将它与原来的不平滑公式混同。

\section{所引用的解析输入}
以下采用本书已经建立的结果，作为本文明确的前提：
\begin{enumerate}
 \item $\zeta$ 的亚纯延拓、函数方程和对数导数的 Dirichlet 级数；
 \item 非平凡零点满足 $0<\Rea\rho<1$，并有 $N(T)=O(T\log(T+2))$，零点计重数；
 \item 对任意固定 $U>0$，存在避零高度 $T_j\to\infty$，使固定竖带
 $-U\le\sigma\le c$ 上 $\zeta'/\zeta(\sigma\pm\ii T_j)
 =O_{U,c}(\log^2(T_j+2))$。
\end{enumerate}
第三项是围道估计输入，不在此重复其局部零点计数与对数导数展开证明。
后文逐步说明如何使用它，而非把围道消失当成无需证明的事实。
这些输入不以 RH 为前提。

\section{先获得绝对收敛的积分表示}
对 $c>1$，由一次平滑的 Perron 核和
$-\zeta'/\zeta(s)=\sum\Lambda(n)n^{-s}$，
\[
 \psi_1(x)=\frac1{2\pi\ii}\int_{c-\ii\infty}^{c+\ii\infty}
 -\frac{\zeta'}{\zeta}(s)\frac{x^{s+1}}{s(s+1)}\dd s.       \tag{1}
\]
合法换序由
$\sum\Lambda(n)n^{-c}<\infty$ 与
$\int_{\mathbb R}|(c+\ii t)(c+1+\ii t)|^{-1}\dd t<\infty$
共同保证。
多出的 $s+1$ 已使原积分线上的反演绝对收敛。

\Needspace{7\baselineskip}
\section{五类留数逐项计算}
被积函数记为 $F_x(s)$。固定 $x>1$，并按重数计零点。
\begin{enumerate}
 \item $s=1$：$\zeta'/\zeta$ 留数为 $-1$，故
 $\operatorname{Res}_{s=1}F_x=x^2/2$。
 \item 非平凡零点 $\rho$：若其重数为 $m_\rho$，
 留数为 $-m_\rho x^{\rho+1}/(\rho(\rho+1))$。
 \item $s=0$：$\zeta(0)\ne0$，核有单极点，留数为
 $-x\zeta'(0)/\zeta(0)=-x\log(2\pi)$。
 \item $s=-1$：此处 $\zeta(-1)=-1/12\ne0$，留数为
 $+\zeta'(-1)/\zeta(-1)$。正号来自 $s=-1$ 处核的留数为 $-1$。
 \item 平凡零点 $s=-2m$（$m\ge1$）：均为单零点，留数为
 $-x^{1-2m}/(2m(2m-1))$。
\end{enumerate}
与未积分的公式相比，$s=-1$ 是新核带来的额外极点，不能遗漏。

\section{围道的极限次序及边积分}
取左边界为 $-U=-(2M+1)$，$M\ge1$。先固定 $M$，沿避零高度令 $T_j\to\infty$。
两条水平边上 $|s(s+1)|\asymp T_j^2$，竖带宽有限，且 $x^{\sigma+1}$ 有固定上界，
故各边积分为 $O_{M,c,x}(\log^2T_j/T_j^2)\to0$。
右竖线积分由绝对收敛趋于式(1)。

再处理无限长的左竖线。由函数方程的对数导数与 Stirling 估计，
在负奇整数线 $\sigma=-U$ 上可得
\[
 \left|\frac{\zeta'}{\zeta}(-U+\ii t)\right|
 \ll\log(U+|t|+2),\qquad U=3,5,\ldots .
\]
这是函数方程估计的直接应用：$1-s$ 位于绝对收敛半平面；
三角因子在负奇整数线上避开其极点，Gamma 对数导数给出对数增长。
在该线上
$|s(s+1)|\gg U^2+t^2$，因而左边积分的模为
\[
 \ll x^{1-U}\int_{\mathbb R}
       \frac{\log(U+|t|+2)}{U^2+t^2}\dd t
 \ll x^{1-U}\frac{\log(U+2)}U\longrightarrow0.
\]
最后一步先令 $t=Uv$，再利用
$\int_{\mathbb R}(1+\log(1+|v|))/(1+v^2)\dd v<\infty$。
这里必须固定 $x>1$；没有证明对 $x\downarrow1$ 一致。

\section{零点和的绝对收敛}
对固定 $x>1$，$|x^{\rho+1}|\le x^2$。
对足够大的 $|\gamma|$，有 $|\rho(\rho+1)|\gg\gamma^2$。
由 Stieltjes 分部积分，
\[
 \sum_{\gamma>T}\gamma^{-2}
 =-\frac{N(T)}{T^2}+2\int_T^\infty\frac{N(u)}{u^3}\dd u
 =O\left(\frac{\log(T+2)}T\right)
\]
（$N$ 仅计正虚部；负虚部由共轭对称处理）。
有限高度的零点项是有限和，且分母不为零，因此全零点级数绝对收敛。
同样的界对任意紧区间 $1<a\le x\le b$ 一致成立。
平凡零点级数则由 $x^{-2m}$ 的几何衰减绝对收敛。

综合留数及边积分，得到
\[
\begin{aligned}
 \psi_1(x)
 ={}&\frac{x^2}{2}
 -\sum_\rho\frac{x^{\rho+1}}{\rho(\rho+1)}
 -x\log(2\pi)+\frac{\zeta'(-1)}{\zeta(-1)}\\
 &-\sum_{m=1}^\infty\frac{x^{1-2m}}{2m(2m-1)},\qquad x>1.
\end{aligned}
\]
这里零点计重数；由于绝对收敛，求和排列不再依赖高度对称规则。

\section{为什么不能随意微分回去}
若逐项微分，分母 $\rho+1$ 被消去，回到约 $x^\rho/\rho$ 的项。
其绝对值上界只有 $O_x(|\gamma|^{-1})$，相应零点级数不绝对收敛。
这里的不绝对收敛不能只由上界推出：固定 $x>1$ 时，
$|x^\rho|=x^\beta\ge1$，而由本书的零点计数渐近式
$N(T)\sim T\log T/(2\pi)$ 及分部求和，$\sum_{\gamma>0}\gamma^{-1}$ 发散。
对充分大的 $|\gamma|$，$|\rho|\asymp|\gamma|$，故求导后的项确实不能绝对求和。
因此本篇的绝对收敛结论不足以支持逐项求导。
恢复 $\psi$ 需要重新处理截断、跳跃和收敛方式；积分获得的优势在求导时可能丢失。

\section*{参考依据与后续核验}
解析输入及特殊值取自本书关于函数方程、零点计数和显式公式的结果；
标准背景可查 E.\ C.\ Titchmarsh, \emph{The Theory of the Riemann Zeta-Function}, 2nd ed., 1986。
本篇直接推导新核、各留数和收敛性，避零高度估计作为引用输入。
作者复审时应重点核对左边界的一致常数与五类留数，不将本初稿视为已经独立重证全部输入。
\end{document}
